\documentclass{beamer}
\usetheme{Madrid}
\usecolortheme{default}
\usepackage{fontspec}
\usepackage{polyglossia}
\setdefaultlanguage{english}
\usepackage{tikz}
\usepackage{hyperref}
\usepackage{graphicx}
\usepackage{booktabs}

% Custom colors
\definecolor{RurDynamixGreen}{RGB}{76, 175, 80}
\definecolor{RurDynamixBlue}{RGB}{33, 150, 243}
\definecolor{DarkGray}{RGB}{64, 64, 64}

\setbeamercolor{primary}{fg=white,bg=RurDynamixGreen}
\setbeamercolor{secondary}{fg=white,bg=RurDynamixBlue}
\setbeamertemplate{frametitle}{
  \begin{beamercolorbox}[rounded=true,shadow=true,leftskip=1cm,rightskip=1cm,wd=\paperwidth,ht=0.8cm,dp=0.4cm]{frametitle}
    \usebeamerfont{frametitle}\insertframetitle
  \end{beamercolorbox}
}

% Title Page Setup
\title{Artificial Basins and Agro-Environmental Dynamics \\ in Bab Berred, Moroccan Rif}
\subtitle{Spatio-Temporal Analysis Using Remote Sensing and GIS (2010-2025)}
\author{Ousmane Samba MAÏGA \\ \small{Supervised by: Prof. Said Moukrim}}
\institute{Rur\_Dynamix Conference 2026 \\ Rural Development and Environmental Dynamics}
\date{\today}

\begin{document}

% Slide 1: Title
\frame{\titlepage}

% Slide 2: Context and Problem Statement
\begin{frame}
\frametitle{Context and Problem Statement}
\begin{itemize}
  \item \textbf{Rapid transformation in Bab Berred}: Cannabis legalization (2021) driving agricultural intensification
  \item \textbf{Adaptation innovation}: Proliferation of artificial basins for irrigation
  \item \textbf{Critical questions}:
  \begin{itemize}
    \item What biophysical factors explain basin location patterns?
    \item What is their actual impact on vegetation cover (NDVI)?
    \item Implications for water sustainability?
  \end{itemize}
  \item \textbf{Urgency}: Climate change + agricultural intensification = hydrological and environmental risks
\end{itemize}
\end{frame}

% Slide 3: Research Objectives
\begin{frame}
\frametitle{Research Objectives}
\textbf{General objective:} Analyze the spatio-temporal dynamics of artificial basins in Bab Berred and assess their impacts on vegetation cover
\\[0.5cm]
\textbf{Specific objectives:}
\begin{enumerate}
  \item Map and quantify basin evolution (2010-2025)
  \item Analyze location factors (elevation, slope, hydrology)
  \item Assess vegetation cover impact (NDVI, 1998-2026)
  \item Identify basin density patterns and cumulative effects
\end{enumerate}
\end{frame}

% Slide 4: Study Area - Bab Berred
\begin{frame}
\frametitle{Study Area: Bab Berred Commune}
\begin{columns}[T]
  \column{0.5\textwidth}
  \textbf{Location:}
  \begin{itemize}
    \item Chefchaouen Province (Central Rif)
    \item Altitude: 1,290 m (mean)
    \item Latitude: 35°00'979"N
    \item Longitude: 004°58'092"W
    \item Average slope: 80% (steep terrain)
  \end{itemize}

  \column{0.5\textwidth}
  \textbf{Population \& Characteristics:}
  \begin{itemize}
    \item 24,742 inhabitants (2024)
    \item Mediterranean mountain climate
    \item Rainfed + irrigation agriculture
    \item Limited water resources
    \item Cannabis cultivation dominant
  \end{itemize}
\end{columns}

\begin{center}
  \small{\textit{Figure 1: Location map of Bab Berred}}
\end{center}
\end{frame}

% Slide 5: Methodology
\begin{frame}
\frametitle{Methodology and Tools}
\textbf{Data sources and tools:}
\begin{itemize}
  \item \textbf{Remote sensing}: Landsat imagery (1998-2026), Google Earth, Sentinel-2
  \item \textbf{GIS}: QGIS, ArcGIS for mapping and spatial analysis
  \item \textbf{Processing}: RStudio, Google Earth Engine (GEE)
  \item \textbf{Indices}: NDVI for vegetation monitoring
\end{itemize}

\textbf{Analytical approach:}
\begin{enumerate}
  \item Manual basin digitization (2010-2025)
  \item Morphometric analysis (surface, volume, depth)
  \item Location analysis (elevation, slope, water network proximity)
  \item Multi-temporal NDVI analysis + precipitation correlation
  \item Spatial density analysis and cumulative effects
\end{enumerate}
\end{frame}

% Slide 6: Basin Growth Results
\begin{frame}
\frametitle{Results: Dramatic Basin Growth in Bab Berred}
\begin{center}
  \textbf{Basin proliferation: Exponential growth 2010-2025}
  \\[0.3cm]
  \small{From dozens to thousands of artificial water storage structures}
\end{center}

\begin{itemize}
  \item \textbf{Acceleration period}: 2013 onwards (legalization in 2021 triggered expansion)
  \item \textbf{Miniaturization trend}: Increasing number of smaller basins
  \item \textbf{Unequal distribution}: Concentration in specific microbasins
  \item \textbf{Spatial anarchy}: Absence of centralized planning or regulation
  \item \textbf{Rapid adoption}: Farmers' response to water scarcity and cannabis demands
\end{itemize}

\begin{center}
  \small{\textit{Figure 2: Cumulative basin count evolution (2010-2025)}}
\end{center}
\end{frame}

% Slide 7: Location Factors
\begin{frame}
\frametitle{Results: Basin Location Factors in Bab Berred}
\textbf{Basin placement responds to biophysical constraints:}
\begin{columns}[T]
  \column{0.5\textwidth}
  \textbf{Elevation:}
  \begin{itemize}
    \item Concentration: 600-1,200 m
    \item Avoids very steep ridges
    \item Avoids flat terrain (poor drainage)
  \end{itemize}

  \textbf{Slope:}
  \begin{itemize}
    \item Preference: 15-35% (water flows)
    \item Terrain adaptation
    \item Efficiency optimization
  \end{itemize}

  \column{0.5\textwidth}
  \textbf{Hydrography:}
  \begin{itemize}
    \item Critical: < 500m to water sources
    \item Water access determines feasibility
    \item Valley bottoms favored
  \end{itemize}

  \textbf{Land cover:}
  \begin{itemize}
    \item Agricultural zones preferred
    \item Forest conversion occurring
    \item Increasing pressure on vegetation
  \end{itemize}
\end{columns}

\begin{center}
  \small{\textit{Figures 3-4: Maps of elevation, slope, and hydrographic distribution}}
\end{center}
\end{frame}

% Slide 8: NDVI and Vegetation Impact
\begin{frame}
\frametitle{Results: Impact on Vegetation Cover (NDVI)}
\textbf{Vegetation dynamics (1998-2026):}
\begin{itemize}
  \item \textbf{Period 1998-2013}: NDVI stable or slightly declining
  \item \textbf{Period 2013-2026}: Significant vegetation improvement
  \item \textbf{Timing}: Rupture coincides with basin expansion onset
  \item \textbf{Correlation}: Strong with irrigation, but also with precipitation cycles
\end{itemize}

\textbf{Interpretation:}
\begin{itemize}
  \item Irrigation via basins → NDVI improvement? (hypothesis)
  \item But also: climate oscillation, natural rainfall variability
  \item Causality not directly proven
  \item Need for fine-scale (plot-level) analysis
\end{itemize}

\begin{center}
  \small{\textit{Figures 5-6: NDVI comparison maps, precipitation correlation}}
\end{center}
\end{frame}

% Slide 9: Basin Morphometrics
\begin{frame}
\frametitle{Results: Basin Morphometric Summary}
\begin{center}
  \small{
  \begin{tabular}{|c|c|c|c|}
    \hline
    \textbf{Parameter} & \textbf{Min} & \textbf{Max} & \textbf{Mean} \\
    \hline
    Surface area (m²) & 50 & 50,000 & 2,500 \\
    Depth (m) & 1 & 8 & 3 \\
    Volume (m³) & 100 & 300,000 & 7,500 \\
    \hline
  \end{tabular}
  }
\end{center}

\textbf{Key trends:}
\begin{itemize}
  \item Progressive miniaturization: older basins > 10,000 m², recent < 1,000 m²
  \item Average basin: 2,500 m² (shrinking trend)
  \item Estimated total volume stored: millions of m³
  \item Implies significant groundwater extraction
\end{itemize}

\begin{center}
  \small{\textit{Figure 7: Distribution of basin dimensions}}
\end{center}
\end{frame}

% Slide 10: Discussion and Sustainability Concerns
\begin{frame}
\frametitle{Discussion: Sustainability Questions}
\textbf{Critical sustainability issues:}
\begin{itemize}
  \item \textbf{Hydrology}: How much water extracted? Groundwater depletion rates?
  \item \textbf{Environmental}: 5,000+ basins = cumulative forest loss, biodiversity impact
  \item \textbf{Social}: Water inequalities, cannabis priority over food crops
  \item \textbf{Climate}: System resilience if precipitation decreases further?
\end{itemize}

\textbf{Knowledge gaps:}
\begin{itemize}
  \item Precise basin-NDVI causality (correlation vs causation)
  \item Hydrologic impact (streamflow changes, aquifer depletion rates)
  \item Long-term viability under climate change
  \item Farmer perspectives and adaptive strategies
\end{itemize}
\end{frame}

% Slide 11: Recommendations
\begin{frame}
\frametitle{Recommendations and Pathways}
\textbf{For research:}
\begin{itemize}
  \item Hydrologic budget studies: basin infiltration, extraction rates
  \item Climate scenario modeling: water availability under future conditions
  \item Participatory research: farmers' views on sustainability
\end{itemize}

\textbf{For territorial management:}
\begin{itemize}
  \item Centralized irrigation planning (avoid spatial chaos)
  \item Remote sensing monitoring system (early warning)
  \item Regulation of basin size/density
  \item Water allocation governance
\end{itemize}

\textbf{For policy:}
\begin{itemize}
  \item Integrate water, agriculture, and climate in rural strategy
  \item Balance cannabis productivity, food security, and environment
  \item Transition toward sustainable practices
\end{itemize}
\end{frame}

% Slide 12: Conclusion
\begin{frame}
\frametitle{Conclusion}
\begin{center}
  \textbf{\Large{Summary and Outlook}}
\end{center}

\textbf{Key messages:}
\begin{itemize}
  \item Bab Berred undergoing rapid transformation driven by cannabis and water storage innovation
  \item Artificial basins = clever local adaptation to water scarcity
  \item But: cumulative impacts poorly quantified → urgent need for monitoring/regulation
  \item NDVI paradox: vegetation greening, but at what cost to groundwater?
\end{itemize}

\textbf{Central challenge:}
\begin{center}
  \textit{Reconciling agricultural productivity, climate resilience, \\ and environmental sustainability in Mediterranean mountain contexts}
\end{center}
\end{frame}

% Final Slide
\begin{frame}[plain]
\begin{center}
  \vfill
  {\Huge \textbf{Thank you for your attention}}
  \\[1cm]
  {\Large Questions?}
  \\[2cm]
  \small{
    Ousmane Samba MAÏGA \\
    Master BioEnv2D \\
    Faculty of Sciences, Rabat
  }
  \vfill
\end{center}
\end{frame}

\end{document}
